\documentclass[12pt,a4paper]{article}
\usepackage[utf8]{inputenc}
\usepackage[T1]{fontenc}
\usepackage{ebgaramond}
\usepackage[french]{babel}
\usepackage[top=35mm,bottom=35mm,left=32mm,right=32mm]{geometry}
\usepackage{microtype}
\usepackage{parskip}
\usepackage{setspace}
\usepackage{xcolor}

\pagestyle{empty}

\begin{document}

% ─── TITRE ────────────────────────────────────────────────────────────────────
\begin{center}
{\scshape\large tête de veau ravigote}

\smallskip
{\itshape Éric Mugnier}

\smallskip
{\small synopsis}
\end{center}

\vspace{10mm}
\noindent\rule{\linewidth}{0.3pt}
\vspace{6mm}

% ─── PITCH ────────────────────────────────────────────────────────────────────
\begin{center}
\begin{minipage}{0.82\linewidth}
\setstretch{1.4}
Quand le père Vidal disparaît dans des circonstances troubles\\
et que des colis macabres commencent à circuler dans la ville,\\
le commandant Beauvais se lance dans une enquête labyrinthique\\
qui le mènera des caves de l'Église à l'ambassade de Namibie.

\medskip
Entouré de Titus Beaugendre, son fidèle équipier,\\
et d'une galerie de personnages hauts en couleur,\\
Beauvais arpente un monde contemporain\\
où l'absurde le dispute à l'horreur,\\
le vulgaire à l'érudit.
\end{minipage}
\end{center}

\vspace{6mm}
\noindent\rule{\linewidth}{0.3pt}
\vspace{8mm}

% ─── DESCRIPTION ──────────────────────────────────────────────────────────────
Un ovni, plutôt qu'un polar. Une logorrhée vengeresse et jubilatoire, écrite en \textit{spontaneous prose}, où l'humour potache s'achoppe à des scènes \textit{gore} à la limite du soutenable~; un roman-fleuve traversé de digressions en forme de cabinet de curiosités~: sur les chats, les dinosaures, l'Église, les pavillons de banlieue, l'unité 731, Pannonica de Koenigswarter, la 'Ndrangheta.

L'\textbf{Horreur} (prononcé à la Marlon Brando dans \textit{Apocalypse Now}) se détache au fil des pages comme le thème central du livre.

Toute forme de morale est étrangère aux protagonistes — au héros-narrateur en premier. Entre grotesque et tragique, les digressions dessinent par intermittence l'image d'une humanité sans Dieu, abominable plus que désespérée. Le ton n'est jamais didactique, jamais emphatique, toujours factuel~— genre «~Alain Decaux raconte~». Un bruit de fond, un acouphène, un filtre optique polarisé~: la tentaculaire progression de la Camorra dans le Naples de l'\textit{Amica Geniale} d'Elena Ferrante, jamais explicite, trop prégnante.

\end{document}
